% !TEX root = ../TCC - CD e IA.tex
\chapter{Conclusões}\label{Conclusões}

% ==========================================================================
% ROTEIRO GERAL DO CAPÍTULO 7
% Sintetizar sem repetir. Responder à pergunta de pesquisa do cap.1 de forma
% direta. Não introduzir resultado novo aqui — só consolidar.
% ==========================================================================

% ROTEIRO (parágrafo de abertura):
% - Retomar a pergunta de pesquisa e afirmar, em 1 parágrafo, como o trabalho
%   a respondeu: a metodologia híbrida, organizada em funil, virou ferramenta
%   de decisão para assessorias.

\section{Síntese das Contribuições}

% ROTEIRO:
% - Recuperar as duas frentes (metodológica e prática) do cap.1.
% - Metodológica: pipeline híbrido integrando 3 técnicas sobre 3 fontes
%   textuais e múltiplos níveis de agregação; refino de tópicos via LLM.
% - Prática: o framework COBRA-RACE, a NSM, o contraste discurso x reação,
%   e as recomendações por perfil.
% - Retomar o achado central: alto engajamento != alta aprovação.
% - Amarrar cada contribuição ao objetivo específico correspondente (1 a 8).

\section{Limitações do Estudo}

% ROTEIRO (ser honesto — a banca valoriza):
% - Dados observacionais, não experimentais (sem A/B): não se estabelece causalidade.
% - Recorte temporal da coleta / poucos pontos para série histórica de CMGR.
% - Transposição de construtos comerciais (COBRAs) ao político — ressalva já
%   feita no cap.3, retomar aqui.
% - Dependência de modelo de sentimento pré-treinado (possíveis vieses em PT-BR,
%   ironia/sarcasmo em comentários políticos).
% - Refino de tópicos via LLM introduz etapa não-determinística.

\section{Sugestões para Trabalhos Futuros}

% ROTEIRO:
% - Teste A/B de formato/tema fechando o loop de aprendizado de growth.
% - Modelo preditivo de engajamento (prever cluster Padrão/Viral/Longo antes de postar).
% - Extensão a outras figuras políticas / outras plataformas (comparabilidade).
% - Análise temporal de retenção com mais coletas.
% - Detecção de ironia/sarcasmo para refinar a análise de sentimento política.

\section{Considerações Finais}

% ROTEIRO:
% - 1 parágrafo de encerramento. Voltar ao "porquê" do trabalho: dar a
%   assessorias uma ferramenta que substitui a intuição por evidência, num
%   ambiente digital decisivo para a democracia. Fechar com a nota de que
%   correlação não implica causalidade — postura que o trabalho manteve.